
% \section{RQ4: What kinds of behaviours to developers overlook in their test suites?}
% \label{sec:rq4}
%\subsubsection{What kinds of behaviours to developers overlook in their test suites?}
% We further examined the untested expected behaviours to understand what kinds of behaviours developers most often miss when constructing their existing test suites, despite these suites having high code coverage. Understanding these patterns can help developers prioritize testing efforts and identify common blind spots in their test-writing strategies. To do this, we used Open Coding (Grounded Theory). The authors iterated through V-UTBs identifying labels that summarized the untested behaviour (e.g., ``Exception'', ``Null handling''). 
% These labels were merged iteratively until saturation was reached. 
% This analysis illustrates that the V-UTBs reported by \toolname represent meaningful categories of behaviours that could be validated by a complete test suite. 
% V-UTBs fall into the following categories: 
\smallskip
\subsubsection{What kinds of behaviours do developers overlook in their test suites?} 
We further examined the UTBs to identify the kinds of behaviours developers frequently miss when constructing their existing test suites, despite these suites having high code coverage. %Understanding these patterns can help developers prioritize testing efforts and identify common blind spots in their test-writing strategies. 
To do this, we used an open coding grounded theory methodology. The authors iterated through a set of valid untested behaviours, identifying labels that summarized each behaviour.
%(e.g., ``Exception'', ``Null handling''). 
These labels were merged iteratively until saturation was reached. % This analysis shows that the validated untested behaviours reported by \toolname fall into meaningful, recurrent categories that could be validated by a more complete test suite.
The most common categories of untested behaviours in developer-written test suites, and an example behaviour of each, includes: 


\begin{itemize}
% \item \rth{it would be nice to get these down to full percentages (e.g., 21\%) to match the rest of the evaluation. but the numbers also only seem to add up to around 81\%?}
    % 20.8\%
    % i think we should remove the percentage here cause i am afraid we would have to justfy the sample size otherwise.
    % ($\approx$21\%)
    \item \textit{Exceptions:} Cases where methods are expected to consistently throw a specific exception under a given condition.
    For example, \textit{``Throws a \texttt{NullPointerException} when the input iterable is null.''}%, or \textit{``Throws an \texttt{IllegalArgumentException} when the provided pattern is invalid.''}  
    % 16.3\%
    %($\approx$16\%)
    \item \textit{Side Effects:} Behaviour describing how a method alters internal state or data beyond what is directly returned. 
    % . Examples include 
    For example, \textit{``Preserves existing entries in the original bag without applying the transformer.''}% or \textit{``Ensuring the context element passed to the method remains unmodified after parsing.''}
    % 13.4\%
    % ($\approx$13\%)
    \item \textit{Handling null/empty inputs:} Explicit documentation of return values or behaviours for null inputs, empty inputs, or absent data.
    % . For instance, 
    For example, \textit{``Returns \texttt{null} when the provided map does not contain the specified key.''}% or \textit{``Return null when the input string is not a valid date representation.''}
    % 9.9\%
    % ($\approx$10\%)
    \item \textit{Construction \& configuration:} Rules describing how objects are initialized or configured. % . For example, 
    For example, \textit{``Constructs a new, empty map with the specified maximum size and load factor.''} %or \textit{``Create the destination directory and any necessary parent directories if they do not already exist.''}
    % 9.9\%
    %  ($\approx$10\%)
    % RTH: i might have collapes this in with `content of returned value', so I'll just leave it off for space.
    % \item \textit{Type of returned object}: Behaviours where the method returns its own instance for chaining or specifies the concrete type or class of the returned object, e.g., 
    % % . Examples: \textit{``Returns the current instance to allow method chaining''} or 
    % \textit{``Produces a new instance of \texttt{HashedMap} that is a separate object from the original.''}  
    % % 9.3\%
    % ($\approx$9\%)
    \item \textit{Content of returned value:} Descriptions of the content or structure of the return value.
    For example, \textit{``Return the previous value of the passed element before it was updated.''}
    % . For example, \textit{``Returns a string containing only valid alphanumeric characters (a--z, A--Z, 0--9) without any special symbols or whitespace.''} or \textit{``Returns an empty Properties object (but not null).''}  
    % 8.3\%
    %  ($\approx$9\%)
    \item \textit{Algorithmic logic:} Descriptions of specific computational rules or constraints. For example, \textit{``Clamps the \texttt{endIndexExclusive} parameter to the array length if it exceeds the array length.''} 
    % 6.4\%
    %($\approx$6\%)
    % RTH: this could be a special case of algorithmic logic (or vice versa) so is left off.
    % \item \textit{Conditional processing}: Behaviours describing how the method handles different input cases, e.g.,
    % % . For example, 
    % \textit{``Use the system's default temporary directory when no parent directory is specified.''} %or \textit{``Return the default value when the value associated with the key cannot be converted to a double.''}
    % 5.7\%
    % ($\approx$6\%)
    % RTH: don't need other since we just said 'main sources'
    % \item \textit{Other}: Covers the remaining UTBs, including interactions with dependencies and immutability/modifiability of returned objects, e.g., \textit{``Applies the provided locale to the date formatting process, using the system's default locale if none is provided.''} %or \textit{``Returns an unmodifiable list that prevents adding or removing elements.''}  
\end{itemize} 


%This categorization captures the range of behaviours identified by \toolname, from exception handling and state changes to configuration rules, algorithmic logic, and API contracts. The \textit{Exception} and \textit{Modification \& Side Effects} categories were among the most frequently observed, suggesting that even mature test suites often overlook critical behaviours related to error handling and state consistency. 
% Our ability to reliably detect these missing test cases could enable us to improve the strength of existing test suites.
% This demonstrates that \toolname identifies meaningful UTBs across important software aspects that remain untested despite being documented.


%\vspace{-0.5ex}
% \begin{rqbox}{What kinds of behaviours to developers overlook in their test suites?}
% {\textbf{Summary:}
% Developer-written test suites fail to validate a wide variety of different behaviours.
% None of the untested behaviour categories are inherently untestable, suggesting existing test suites could be strengthened by including them.
% }
% \end{rqbox}
%\vspace{-3ex}

% \begin{tcolorbox}[colback=gray!10!white, colframe=black, boxrule=0.5pt]
% \textbf{Answer to RQ2:}
% Developer-written test suites miss a broad set of expected behaviours, even for high-coverage test suites.
% By identifying these missed behaviours, \toolname could help developers more comprehensively validate their systems.
% \end{tcolorbox}

% \subsection{RQ3: Can \toolname generate valid test cases for untested behaviours?}

% \toolname seeks to generate new test cases that can be used to strengthen an existing suite by adding tests for untested behaviours. 
% To evaluate this we conducted a two-phase evaluation.
% The first phase evaluates whether \toolname can actually generate test cases for untested behaviours.
% As with RQ1 we lack an oracle to compare these test cases against (and they do not exist in the existing suite), so this phase of the evaluation focuses on the proportion of untested behaviours \toolname can generate test cases for.
% The benefit of this design is that we can evaluate the test generation across all 3,083 untested behaviours uncovered in RQ1 to ensure we can generate test cases across a variety of projects and behaviours.
% The drawback of this design is that this evaluation does not tell us if the generated test cases would actually useful. 

% To address this drawback, we performed a second phase to manually examine 100 generated test cases from each project. This manual analysis sought to answer whether the generated test cases could strengthen the test suite by evaluating untested behaviours.
% The benefit of this manual analysis is that it helped us understand whether the generated tests were actually useful.
% The drawback of this analysis is that it was a manual effort performed by a single author.

% \paragraph{Can \toolname generate test cases for untested behaviours?}

% Table~\ref{tab:total_test_methods} details the pool of untested behaviours and the test cases that \toolname generates for them.
% This shows that the approach is able to generate test cases for the vast majority of untested behaviours.
% Approximately 6\% of the generated test cases raised assertion errors when they were executed; these are described further in the second phase of this evaluation.

% The newly-generated test cases did not meaningfully increase coverage, which is expected given the already high line and branch coverage of the existing test suites.
% At the same time though, the non-failing generated test cases significantly decreased the proportion of untested behaviours in the test suites from 3,083 to 645. To ensure these generated tests were actually valid, we next manually investigated a subset of these generated tests, along with the generated tests that raised assertion errors.

% \begin{table}[h]
% \centering
% \caption{Results for the test case generation. \toolname was able to generate test cases for 85\% of UTBs; only 6\% of these generated tests cases failed at runtime due to assertion errors. As expected, the coverage gain from these new test cases was minimal.}
% \label{tab:total_test_methods}
% \begin{tabular}{l S[table-format=4.0] ccc}
% \toprule
% \textbf{Project} & \textbf{UTB} & \textbf{Generated Unit Tests} & \textbf{Assertion Failures} & \textbf{$\Delta$ Coverage} \\
% \midrule
% commons-collections & 1,102 & \hphantom{1,}979\space($88.8\%$) & 44 & 0.50\% \\
% commons-lang        & 1,403 & 1,160\space($82.7\%$) & 88 & 0.30\% \\
% jsoup               & 282 & \hphantom{1,}222\space($78.7\%$) & 19 & 0.10\%\\
% plexus-utils        & 161 & \hphantom{1,}140\space($87.0\%$) & 11 & 1.90\%\\
% commons-cli         & 135 & \hphantom{1,}114\space($84.4\%$) & 15 & 0.07\% \\
% \bottomrule
% \end{tabular}
% \end{table}
% \noindent

% % We additionally manually examined an additional 177 test cases that 

% % Most generated test cases compiled successfully. A small fraction failed at runtime due to assertion errors, and we reported separately on why they failed in section~\ref{par:assertion-failures}.

% \paragraph{Are the generated test cases correct?}

% To evaluate correctness, one author manually reviewed a random sample of 100 generated test cases for each project. Each case took approximately 2–3 minutes to validate. For each generated test case we performed a two-step process:

% \begin{enumerate}
%     \item \textit{Correct method and untested behaviour.} We ensured that the generated test actually invoked the intended method and untested behaviour.
%     \item \textit{Correct assertions.} We also examined the generated assertions to ensure they tested the expected behaviour rather than simply repeating method return values. We also executed the test cases to confirm their runtime behaviour. 
% \end{enumerate}

% If both of these steps yielded a positive result, the test case was considered \emph{valid}, otherwise we considered it \emph{invalid}; the results of this manual process is shown in Table~\ref{tab:manual_validation_generated_tests}.
% Most ($\approx$95\%) of the generated test cases were valid. Invalid cases typically resulted from misunderstandings of documentation, incorrect assumptions about default parameter values, or incomplete execution context. 
% In particular, \toolname sometimes lacked sufficient background information about how a method should be set up or invoked. For example, some methods required the invocation of the parent class method or values, external dependencies, or specific object states before invocation, but this context was not always inferable from the extracted method context alone.

% \begin{table}[h]
% \centering
% \caption{Results of the manual generated test validation. $\approx$95\% of the generated tests were valid.}
% \label{tab:manual_validation_generated_tests}
% \begin{tabular}{lccc}
% \toprule
% \textbf{Project} & \textbf{Samples} & \textbf{Valid} & \textbf{Invalid} \\
% \midrule
% commons-collections & 100 & 98 & 2 \\
% commons-lang        & 100 & 95 & 5 \\
% jsoup               & 100 & 94 & 6 \\
% plexus-utils        & 100 & 95 & 5 \\
% commons-cli         & 100 & 91 & 9 \\
% \bottomrule
% \end{tabular}
% \end{table}

% \paragraph{What caused the assertion failures?}
% \label{par:assertion-failures}

% In addition to the manual validation of the sampled test cases, we also analyzed cases where tests compiled but failed at runtime due to assertion failures. These failures did not always indicate incorrect test generation but rather highlighted misalignment between the method's documentation and its implementation. We grouped the observed failures into the following categories:
% \begin{itemize}
% \item \textit{LLM hallucinations (91\%):} The majority of failures were a result of the LLM generating tests for behaviours not supported by the method. This category included (i) \emph{invalid inputs or assumptions}, where the test relied on unsupported defaults or parameter states and incorrect outputs, and (ii) \emph{over-generalization}, where the LLM invented entirely new behaviour that were absent from both documentation and implementation. 

%     \item \textit{Context or dependency gaps (6\%):} These failures occurred when the tool did not have enough necessary context in the prompt, such as properties from parent classes, external state, or dependent objects.
%     %In these cases, the test logic was otherwise correct, but the incomplete context led the LLM to hallucinate about properties or object state. 

%     \item \textit{Documentation/implementation mismatches (3\%):} These failures occurred when the documentation described a behaviour that diverged from the actual implementation. 
%     % For example, in the \textit{commons-lang} project, the \texttt{StrSubstitutor} constructor was documented to throw an \texttt{IllegalArgumentException} if the prefix was null, but the implementation instead threw a \texttt{NullPointerException}. Generated tests that followed the documentation failed at runtime due to this mismatch (\autoref{mismatches}). Such cases reveal real inconsistencies between specification and code. 
% \end{itemize} 

% One such example of \textit{Documentation/implementation mismatches} is the \texttt{StrSubstitutor} constructor from \texttt{commons-lang} where the documentation states:
% \begin{quote}
% \small
% \texttt{@throws IllegalArgumentException if the prefix or suffix is null}
% \end{quote}
% Based on the behaviour inferred from this line of documentation, \toolname generates the test case shown in Figure~\ref{mismatches}. 
% This generated test fails at runtime because the implementation for \texttt{StrSubstitutor} conflicts with its documentation: the code throws a \texttt{NullPointerException} when the prefix is null, not an \texttt{IllegalArgumentException}.
% Such cases highlight both the importance of aligning the documentation with the implementation it describes, but also in perform fully quantitative studies that do involve any manual analysis that examines the nature of the test cases that are failing.

% \begin{figure}[h]
% % \begin{tcolorbox}[colback=gray!5, colframe=black!60, title=Mismatch between documentation and implementation.]

% \begin{lstlisting}[language=Java]
% @Test
% void testConstructorNullPrefix() {
%     final Map<String, String> map = new HashMap<>();
%     map.put("name", "commons");

%     // Test with null prefix
%     IllegalArgumentException exception = assertThrows(IllegalArgumentException.class,
%         () -> new StrSubstitutor(map, null, ">", '!')
%     );
% }
% \end{lstlisting}
% % \end{tcolorbox}
% \caption{Generated test case for checking \texttt{IllegalArgumentException} exception for \texttt{StrSubstitutor} constructor based on the documentation.}
% \label{mismatches}
% \end{figure}

% Overall, the generated test cases we examined throughout this evaluation strengthened the existing developer-written suites in meaningful ways by covering new behaviours not previously validated and by providing new forms of documentation through executable examples. 

% \begin{tcolorbox}[colback=gray!10!white, colframe=black, boxrule=0.5pt]\textbf{Answer to RQ3:}
% \toolname generates test cases for $85\%$ of the untested behaviours it is able to identify. Manual validation of generated tests across five projects showed the tests were valid $\approx$95\% of the time. 
% % \rth{need this to be consistent with whatever we choose for the tab 4 caption}
% %in between 91–98\%, with the main causes of invalidity being LLM hallucinations, missing context, or ambiguity in the source documentation and actual implementation.
% \end{tcolorbox}





% RTH: if i had to drop a RQ, the ablation is the one I'd leave out
\endinput 

\subsection{RQ4: What is contribution do the major components make to the overall effectiveness of the approach?}

%Our overall approach is composed of two main components: the \textit{untested behaviour identifier} and the \textit{test generator}. Since these components serve different purposes, we study them separately. We designed one ablation to better understand how the design decisions behind the two components in Figure~\ref{fig:tool} contributed to its overall effectiveness. 

We designed an ablation study to examine how the different design decisions behind \toolname{}'s overall workflow in Figure~\ref{fig:tool} contributed to its overall effectiveness.

%understand how much each step of the untested behaviour identifier contributes to identifying untested behaviours, and another ablation to evaluate which elements of the test generator are essential for producing valid test cases. 


%\subsubsection{Untested behaviour identifier variants}

To understand the contribution of each system component in the UTB Identifier, we perform an ablation study to evaluate the importance of splitting the behaviour identification step for a method from the identifying untested behaviours with the test cases.

Unlike RQ1, where we sampled individual UTBs, here we focus on a set of methods from the validation set. For each method, we compare our full pipeline with two ablated variants. Since the number of predicted UTBs can vary across settings, we report \emph{precision}, the number of \emph{invalid UTBs}, and \emph{Already Tested}. Here, invalid refers to hallucinated behaviours not described in the method’s documentation, and these metrics directly reflect how effectively each component constrains the generation of invalid untested behaviour. Invalid UTBs therefore capture not only failures in mapping behaviours to tests, but also cases where the model infers expected behaviours that are not supported by the method’s documentation.

\paragraph{Baseline (A0: \toolname)} 
The full pipeline uses both the multi-step reasoning and additional context, like the call graph path, and the target method's implementation. 
In this way the tool reasons about the behaviours of the method under test and the test cases in separate steps.
% This allows the LLM to fully reason about which behaviours are already tested and which are genuinely untested.

\paragraph{Ablation 1 (A1: One-shot UTB prediction).}

In this setting, we remove the multi-step structure. We provide the LLM with the method and its documentation along with all related test cases (within the same 2-hop call graph distance). The LLM is asked to list the UTBs for that method directly.
%
This leads to a large drop in precision across all repositories. The LLM often predicts more UTBs, many of which are invalid. We found that it frequently focuses on describing what is missing from the test cases, rather than aligning with the documented behaviours. Additionally, without separating behavioural extraction and test mapping, the model often fails to recognize what has already been tested. For instance, in \texttt{commons-lang}, invalid UTBs jumped from 7 to 93. This suggests that separating behavioural extraction from test mapping is essential for constraining the LLM to documented requirements.

\paragraph{Ablation 2 (A2: Removing additional context).}

In this setting, we retain the full multi-step structure of our pipeline but remove the additional context, like the call graph path that connects the test cases to the target method and its implementation. 
The LLM is provided with the list of behaviours and the related test case, but without any context showing how the tests might reach the target functionality or the implementation of the target method.
%
Compared to Ablation 1, this configuration results in fewer invalid behaviours and achieves higher precision. However, it introduces a significant increase in Already Tested false positives, which often match or exceed the counts seen in A1.
This suggests, without the call path or the method implementation, the LLM struggles with multi-hop reasoning and often fails to detect that certain behaviours are being tested. As a result, several tested behaviours are incorrectly flagged as untested. While A2 shows variable improvement over A1 (5-36 percentage points depending on project complexity), it consistently underperforms the full pipeline by 15-25 percentage. Notably, A2 consistently outperforms A1 despite both having significant limitations. This demonstrates that the multi-step decomposition provides value even without call graph context, though both components are necessary for optimal performance.
% As a result, although precision improves relative to A1, it remains substantially lower than that of the full pipeline (A0).

Table~\ref{tab:ablation-results} shows the number of invalid UTBs, Already Tested count, and precision across both ablation settings, along with the baseline. These results show that while the step-by-step breakdown helps constrain the task, the call graph context is essential for enabling the LLM to reason effectively across multiple layers of indirection. Removing this context consistently increases invalid (hallucinated) behaviours, even when the overall task structure is preserved.


\begin{table}[h]
\centering
\caption{Ablation results across three configurations (A0: \toolname, A1: one-shot prompt, A2: no call graph).}
\label{tab:ablation-results}
\begin{tabular}{lccccccccc}
\toprule
\multirow{2}{*}{\textbf{Project}} 
& \multicolumn{3}{c}{\textbf{Invalid UTBs}} 
& \multicolumn{3}{c}{\textbf{Already Tested}} 
& \multicolumn{3}{c}{\textbf{Precision (\%)}} \\ \cline{2-10} 
& A0 & A1 & A2 & A0 & A1 & A2 & A0 & A1 & A2\\ \hline
commons-collections & 6  & 27 & 10 & 46 & 71 & 67 & \textbf{60.3} & 35.1 & 40.1 \\
commons-lang        & 7  & 93 & 14 & 16 & 50 & 29 & \textbf{84.7} & 38.3 & 63.1 \\
jsoup               & 3  & 15 & \hphantom{1}2 & \hphantom{1}4 & 28 & \hphantom{1}6 & \textbf{85.6} & 34.4 & 70.2 \\
plexus-utils        & 1  & 36 & \hphantom{1}2 & \hphantom{1}2 & 22 & 17 & \textbf{93.2} & 40.3 & 56.7 \\
commons-cli         & 2  & 20 & \hphantom{1}1 & \hphantom{1}4 & 21 & 16 & \textbf{85.3} & 36.7 & 42.1 \\ \bottomrule
\end{tabular}
\end{table}

These results align with our earlier finding (RQ3) that incomplete test identification is the primary source of false positives. The A2 configuration, which removes call graph context, shows a 2-4x increase in `Already Tested' false positives, confirming that test identification accuracy drives overall precision.

\begin{tcolorbox}[colback=gray!10!white, colframe=black, boxrule=0.5pt]\textbf{RQ4: Ablation Summary for UTB Identification:}
The full pipeline (A0) achieves the highest precision, showing the benefit of processing the task step by step while using the call graph context. 
These results confirm that both stepwise processing and the call graph context, including method implementation, are essential to reliably identify untested expected behaviours.
\end{tcolorbox}
%Removing the step-by-step approach (A1) greatly increases invalid UTBs and lowers precision, while removing the call graph paths and the target method's implementation (A2) improves slightly over A1 but still underperforms the full pipeline. 

% \subsubsection{Test generator variants} 
 
%  We next ablated the test generator. The baseline pipeline (\textit{G0}) integrates project-level context (available method signatures) and a \textit{Validator} component that iteratively repairs generated code. We compared against one ablated variant: (\textit{G1}) without the \textit{Validator} component. We did not remove the project-level context, as it provides essential information about available method signatures and class structure, which the LLM requires to generate valid test cases. Table \ref{tab:generator_ablation} shows the number of successfully generated unit tests (compilation error-free) and their proportion relative to total UTBs. We only count test cases that compiled successfully and we excluded generated test cases with compilation errors.

% \begin{table}[h]
% \centering
% \caption{Ablation results across both configurations of the test generator for successfully generated test cases.}
% \label{tab:generator_ablation}
% \begin{tabular}{l S[table-format=3.0] c c c }
% \toprule
% \textbf{Project}& \textbf{UTB} & \textbf{G0: \toolname} & \textbf{G1: w/o Validator} \\
% \midrule
% commons-collections &1,102 & \textbf{\hphantom{1,}979 (88.8\%)} & 497 (45.2\%)\\
% commons-lang        & 1,403 & \textbf{1,160 (82.7\%)} & 605 (43.1\%) \\
% Jsoup               & 282  & \textbf{\hphantom{1,}222 (78.7\%)} & 142 (50.3\%)  \\
% Plexus-Utils        & 161  & \textbf{\hphantom{1,}140 (87.0\%)} & \hphantom{1}78 (48.7\%)  \\
% commons-cli         &135  & \textbf{\hphantom{1,}114 (84.4\%)} & \hphantom{1,}60 (44.8\%)\\
% \bottomrule
% \end{tabular}
% \end{table}

%\begin{tcolorbox}[colback=gray!10!white, colframe=black, boxrule=0.5pt]\textbf{RQ3: Ablation Summary for Test Generation:}
%Without the \textit{Validator} (\textit{G1}), only 45\% 
%of the generated tests successfully compiled, which is a substantial decline from the 85\% for the baseline. This confirms that the \textit{Validator} is crucial to repair the semantic and syntactic errors of the generated test cases for untested behaviours.
%\end{tcolorbox}

% \begin{table}[ht]
% \centering
% \caption{Ablation results across two configurations.}
% \label{tab:ablation-results}
% \begin{tabular}{|l|c|c|c|c|}
% \hline
% \textbf{Repository} & \textbf{Ablation} & \textbf{Extracted UTBs} & \textbf{Invalid UTBs} & \textbf{Precision (\%)} \\
% \hline
% commons-collections & A1 & 174 & 27 & 35.14 \\
% commons-collections & A2 & 145 & 10 & 40.14 \\
% commons-lang        & A1 & 170 & 93 & 38.34 \\
% commons-lang        & A2 & 152 & 14 & 63.08 \\
% jsoup               & A1 & 61  & 15 & 34.38 \\
% jsoup               & A2 & 39  & 2  & 70.21 \\
% plexus-utils        & A1 & 76  & 36 & 40.30 \\
% plexus-utils        & A2 & 64  & 2  & 56.66 \\
% commons-cli         & A1 & 57  & 20 & 36.73 \\
% commons-cli         & A2 & 39  & 1  & 42.07 \\ \hline
% \end{tabular}
% \end{table}

% \begin{table}[ht]
% \centering
% \caption{Ablation results across two configurations.}
% \label{tab:ablation-results}
% \begin{tabular}{|l|c|c|c|c|}
% \hline
% \textbf{Repository} & \textbf{Ablation} & \textbf{Invalid UTBs} & \textbf{Precision (\%)} \\
% \hline
% commons-collections & A0 & 6 & 60.3\% \\

% commons-collections & A1 & 27 & 35.14 \\
% commons-collections & A2 &  10 & 40.14 \\

% commons-lang        & A0 &  7 & 84.7\% \\

% commons-lang        & A1 &  93 & 38.34 \\
% commons-lang        & A2 &  14 & 63.08 \\

% jsoup               & A0 & 2 & 85.6 \\

% jsoup               & A1 & 15 & 34.38 \\
% jsoup               & A2 &  2  & 70.21 \\

% plexus-utils        & A0 &  1 & 93.1\%\\

% plexus-utils        & A1 &  36 & 40.30 \\
% plexus-utils        & A2 &  2  & 56.66 \\

% commons-cli         & A0 &  2 & 85.3\% \\

% commons-cli         & A1 &  20 & 36.73 \\
% commons-cli         & A2 &  1  & 42.07 \\ \hline
% \end{tabular}
% \end{table}
